% Page de garde — Mémoire de Master, UNIKIN
% Compilation : xelatex page_de_garde.tex
\documentclass[12pt,letterpaper]{report}
\usepackage{fontspec}
\usepackage[french]{babel}
\usepackage[letterpaper,top=25mm,bottom=25mm,left=30mm,right=25mm]{geometry}
\usepackage{graphicx}
\usepackage{setspace}
\begin{document}
\begin{titlepage}
\begin{center}
\singlespacing
\includegraphics[width=2.4cm]{blason_unikin.png}\\[0.4cm]
{\large\bfseries UNIVERSITÉ DE KINSHASA}\\[0.35cm]
{\bfseries FACULTÉ DES SCIENCES ET TECHNOLOGIES}\\[0.12cm]
{\bfseries MENTION DE MATHÉMATIQUES, STATISTIQUE ET INFORMATIQUE}\\[0.12cm]
B.P. 190 KINSHASA XI\\[1.6cm]
\rule{\linewidth}{0.5pt}\\[0.3cm]
{\Large\bfseries Détection et localisation des segments générés dans des textes hybrides francophones : la variance intra-document des traits linguistiques comme signal\par}
\vspace{0.2cm}
\rule{\linewidth}{0.5pt}\\[0.5cm]
{\small Mémoire présenté et défendu en vue de l’obtention du Grade de Master en Sciences et Technologies}\\[0.5cm]
{\small Présenté par}\\[0.15cm]
{\large\bfseries MASIKOTI FUTU Jérémie}\\[0.4cm]
Filière : \textbf{Informatique}\\[0.1cm]
Option : \textbf{Intelligence artificielle et Sciences des données}\\[0.6cm]
Directeur : {\bfseries Prof. Dr Selain KASEREKA KABUNGA}
\vfill
{\bfseries ANNÉE ACADÉMIQUE 2025--2026}
\end{center}
\end{titlepage}
\end{document}
